% Section: Conclusion
\section{Conclusion}
\label{sec:conclusion}

We introduced \textbf{HiddenCourtesyMerge-Sim}, a controlled benchmark for
belief-based planning under hidden driver courtesy in highway merge scenarios.
The benchmark provides a 3{,}000-episode dataset with ground-truth binary
latent labels, a calibrated product-of-Gaussians observation model
(98.3\% held-out step-level accuracy, NLL 0.168, $\lambda{=}0.08$ selected
by grid search), an abstracted hidden-parameter POMDP interface with Bayesian
belief update, and a five-policy matched evaluation with paired statistical
evaluation infrastructure
(20-test joint Holm--Bonferroni, paired Cohen's $d$, post-hoc power analysis).
Oracle access to the hidden type significantly reduces collisions
(10.3\% vs.\ 18.3\%, $p_\text{adj}{=}0.005$, $|h|{=}0.230$) but a Bayesian
belief policy does not close this gap at $n{=}300$ ($p_\text{adj}{=}0.33$,
achieved power 19.5\%); a belief-gain sweep localises the bottleneck in
early-window observation quality, not action-table expressiveness.
A POMCP solver achieves 5.0\% collision and reward 57.36, significantly
exceeding all heuristic baselines, including the oracle-heuristic baseline
($p_\text{adj}{=}0.014$, $|h|{=}0.204$, $d{=}0.221$), indicating within this
benchmark that UCT planning over the reward function captures
safety-relevant action sequences not represented by the heuristic action table.
\textit{The benchmark's primary value is providing the matched experimental
infrastructure and mechanistic diagnostics---$\lambda$-ablation,
early-window accuracy, belief-gain sweep, reward sensitivity, rollout
isolation, covariance ablation, and paired power analysis---needed to clearly
distinguish what online belief estimation and principled planning each
contribute to merge safety.}
An anonymized review repository will be provided during submission; the
de-anonymized dataset, code, and figures will be released upon acceptance.
